\clearpage
\VerdierPhysicalPageAnchor{0134}%
\begingroup
\small
\setlength{\emergencystretch}{1.5em}%
\setlength{\abovedisplayskip}{1pt plus 1pt minus 1pt}%
\setlength{\belowdisplayskip}{1pt plus 1pt minus 1pt}%

\VerdierSemanticAnchor{ch2-2-1-15}
\noindent\textbf{2.1.15.}\quad
\VerdierSemanticAnchor{ch2-2-1-15-associated-system}%
$\mathcal D'$ を三角圏 $\mathcal D$ の充満三角部分圏とする。乗法的系
$S(\mathcal D')$ (2.1.7) を、\emph{部分圏 $\mathcal D'$ に付随する
乗法的系}と呼ぶ。

\VerdierSemanticAnchor{ch2-2-1-16}
\smallskip
\noindent\textbf{命題 2.1.16.}\quad
\begingroup\itshape
\VerdierSemanticAnchor{ch2-2-1-16-setup}%
$F\colon\mathcal D_1\to\mathcal D_2$ を二つの三角圏の間の完全関手とする
(1.1.3)。$S(F)$ を、$F$ により $\mathcal D_2$ の同型へ送られる
$\mathcal D_1$ の射の集合とする。
\VerdierSemanticAnchor{ch2-2-1-16-kernel-subcategory}%
$\mathcal D'(F)$ を、$F$ により $\mathcal D_2$ の零対象へ送られる
$\mathcal D_1$ の対象で定まる $\mathcal D_1$ の充満部分圏とする。

\VerdierSemanticAnchor{ch2-2-1-16-a}
\par\noindent\hspace*{2em}a)\enspace 部分圏 $\mathcal D'(F)$ は厳密充満
三角部分圏（$\mathcal D'(F)$ の対象と同型な $\mathcal D_1$ の任意の
対象は $\mathcal D'(F)$ の対象である）であり、飽和である。

\VerdierSemanticAnchor{ch2-2-1-16-b}
\par\noindent\hspace*{2em}b)\enspace 集合 $S(F)$ は $\mathcal D_1$ の
飽和乗法的系である。

\VerdierSemanticAnchor{ch2-2-1-16-c}
\par\noindent\hspace*{2em}c)\enspace 集合 $S(F)$ は部分圏
$\mathcal D'(F)$ に付随する乗法的系である (2.1.15)。
\par\endgroup

\VerdierSemanticAnchor{ch2-2-1-16-proof-omitted}
\noindent 証明は容易なので読者に委ねる。同じ方法で次の命題も証明できる：

\VerdierSemanticAnchor{ch2-2-1-17}
\smallskip
\noindent\textbf{命題 2.1.17.}\quad
\begingroup\itshape
\VerdierSemanticAnchor{ch2-2-1-17-setup}%
$H\colon\mathcal D\to\mathcal A$ を三角圏からアーベル圏への
コホモロジー的関手とする。$S(H)$ を、$\mathcal D$ の射で、そのすべての
シフトとともに $\mathcal A$ の同型へ送られるものの集合とする。
\VerdierSemanticAnchor{ch2-2-1-17-kernel-subcategory}%
$\mathcal D'(H)$ を、$\mathcal D$ の対象で、そのすべてのシフトとともに
$\mathcal A$ の零対象へ送られるものによって定まる $\mathcal D$ の充満
部分圏とする。

\VerdierSemanticAnchor{ch2-2-1-17-a}
\par\noindent\hspace*{2em}a)\enspace 部分圏 $\mathcal D'(H)$ は厳密充満
三角部分圏であり、飽和である。

\VerdierSemanticAnchor{ch2-2-1-17-b}
\par\noindent\hspace*{2em}b)\enspace 集合 $S(H)$ は $\mathcal D$ の飽和
乗法的系である。

\VerdierSemanticAnchor{ch2-2-1-17-c}
\par\noindent\hspace*{2em}c)\enspace 集合 $S(H)$ は部分圏
$\mathcal D'(H)$ に付随する乗法的系である (2.1.15)。
\par\endgroup

\VerdierSemanticAnchor{ch2-2-2}
\medskip
\noindent\textbf{2.2.\quad 局所化圏の構成。}

\VerdierSemanticAnchor{ch2-2-2-1}
\smallskip
\noindent\textbf{2.2.1.}\quad
\VerdierSemanticAnchor{ch2-2-2-1-slice-categories}%
$\mathcal C$ を圏（加法的であるとは限らない）、$S$ を $\mathcal C$ の
乗法的系とする。$\mathcal C$ の各対象 $X$ に対し、$S/X$（それぞれ
$X\backslash S$）を、対象が終域 $X$（それぞれ始域 $X$）をもつ $S$
の射であり、
\VerdierSemanticAnchor{ch2-2-2-1-morphism-diagrams}%
射が次の可換図式であるような圏とする：

\VerdierSemanticAnchor{ch2-2-2-1-s-over-x-diagram}
\begin{center}
  \begin{tikzpicture}[>=stealth,line width=.38pt,scale=.90,transform shape]
    \node (Y)  at (-1.55,1.05) {$Y$};
    \node (Yp) at ( 1.55,1.05) {$Y'$};
    \node (X)  at ( 0,-.95) {$X$};
    \draw[->] (Y) -- node[above=1pt] {$u$} (Yp);
    \draw[->] (Y) -- node[left=2pt] {$s$} (X);
    \draw[->] (Yp) -- node[right=2pt] {$t$} (X);
    \node at (3.15,.05) {$s,t\in S$};
  \end{tikzpicture}
  \quad ,
\end{center}

\VerdierSemanticAnchor{ch2-2-2-1-p135-boundary}
\endgroup
